I'll write the assigned LaTeX section using only the shared facts provided.

`\section{Exploratory Data Analysis and Data Pipeline}`
`\label{sec:methods_data}`

`\subsection{Exploratory data analysis}`
The training corpus is SIPaKMeD~`\citep{plissiti2018sipakmed}`, a collection of 5{,}015 Papanicolaou-stained conventional-smear images annotated into five cytomorphology classes: Superficial--Intermediate, Parabasal, Metaplastic, Koilocytotic, and Dyskeratotic. Three properties of this corpus shape the entire downstream pipeline, and we summarize them visually in Figure~`\ref{fig:eda}`.

First, the class distribution is only mildly imbalanced. The per-class counts---Dyskeratotic 1{,}036, Koilocytotic 1{,}063, Metaplastic 1{,}064, Parabasal 895, and Superficial--Intermediate 957---yield a majority-to-minority ratio of 1.19. This near-uniform support means that the dominant modelling challenge is morphological discriminability among adjacent cytological categories rather than gross frequency skew; nevertheless, we retain a class-weighted training objective (Section on training) as a conservative safeguard against the residual imbalance.

Second, the imagery is strongly heterogeneous in spatial extent. Image widths span 62--2{,}048~px (mean 614.7) and heights span 62--1{,}536~px (mean 485.4), a roughly thirty-fold range in each dimension. Such variability precludes feeding native resolutions to a convolutional backbone and motivates an explicit, fixed-size resizing stage that places every sample on a common geometric footing.

Third, the dataset exhibits a dual granularity: 4{,}049 of the images are tightly cropped single cells, while the remaining 966 are cluster field-of-view (FoV) images containing multiple overlapping cells. This mixture is biologically faithful---routine smears contain both isolated cells and dense clusters---but it has an important methodological consequence. Cells originating from the same FoV are visually correlated, so a naive per-image split can place near-duplicate views on both sides of a train/validation partition. We flag this here because it directly conditions our later interpretation of the optimistic internal accuracy.

`\subsection{Preprocessing}`
All images pass through a single deterministic preprocessing chain at both training and evaluation time. Each file is decoded and converted to three-channel RGB, then resized to $380\times380$~px, the native input resolution of EfficientNet-B4~`\citep{tan2019efficientnet}`. Crucially, we adopt this same $380\times380$ geometry for \emph{both} the B4 teacher and the B0 student, so that distillation transfers knowledge between networks that observe pixel-identical inputs and the student is never disadvantaged by a resolution mismatch. Pixel intensities are finally standardized with the ImageNet channel statistics---means $(0.485, 0.456, 0.406)$ and standard deviations $(0.229, 0.224, 0.225)$---consistent with the ImageNet-pretrained weights~`\citep{deng2009imagenet}` that initialize both backbones. Fixing resolution and normalization across teacher, student, and all four datasets ensures that any cross-dataset performance gap reflects genuine domain shift rather than incidental differences in the input pipeline.

`\subsection{Data augmentation}`
Stochastic augmentation is applied \emph{exclusively} during training; the validation and all external evaluations use only the deterministic resize-and-normalize chain above, so that reported generalization is never inflated by test-time perturbation. The training-time transforms are deliberately restricted to operations that preserve cytological meaning. We apply random horizontal and vertical flips (each with probability 0.5), random rotation within $\pm72^{\circ}$, and a random affine scaling (zoom) between $0.9$ and $1.1$. These choices exploit a defining symmetry of cytology: a cell has no canonical orientation on the slide, so reflections and rotations generate label-preserving views that encourage rotation- and flip-invariant features rather than memorized pose. The modest scaling band emulates benign variation in cell size and magnification without distorting the nucleus-to-cytoplasm relationships that underpin the diagnostic classes. We deliberately omit colour, contrast, and elastic distortions, which could corrupt the Papanicolaou staining cues and chromatin texture that the classes depend on. This augmentation regime enlarges the effective sample diversity while keeping every synthetic view cytologically plausible.